% =====================================================================
%  CHAPITRES 3 ET 4 — Les techniques d'entraînement + les outils
% =====================================================================
\chapter{Les huit techniques qui font parler}

Une méthode de langue est une boîte à outils. Voici les huit outils qui produisent
l'essentiel des résultats. Chacun a son protocole précis --- appliqué sérieusement, pas
approximativement. Maîtrisez-les tous dans le premier mois : ensuite, votre année n'est plus
qu'une question de contenu.

\section{Technique 1 --- Le \emph{shadowing} (l'ombre)}
\textbf{Ce que c'est :} écouter un audio et le répéter \emph{presque en même temps} que le
locuteur, comme une ombre qui suit, avec un décalage d'une demi-seconde. C'est l'exercice le
plus puissant pour la prononciation, le rythme, l'intonation et la vitesse de parole.

\begin{tableau}[Protocole du shadowing --- 7 étapes, 15 minutes]
\begin{tabularx}{\linewidth}{@{}p{1.2cm}p{4.2cm}X@{}}
\entete{Étape} & \entete{Action} & \entete{Détail} \tabularnewline
\hline
1 & Choisir un audio court & 30 à 90 secondes, niveau compréhensible à 90\,\%, voix que vous aimez. Toujours le même pendant 3 jours. \tabularnewline
2 & Écouter sans texte & 2 fois. Comprendre le sens global. \tabularnewline
3 & Lire le texte & 1 fois, à voix haute lentement, en cherchant les mots inconnus. \tabularnewline
4 & Écouter en lisant & 2 fois, en suivant du doigt. Repérer les liaisons et les mots « mangés ». \tabularnewline
5 & \textbf{Shadower avec le texte} & 3 fois : vous parlez en même temps que l'audio, les yeux sur le texte. Imiter \emph{tout} : mélodie, pauses, énergie. \tabularnewline
6 & \textbf{Shadower sans le texte} & 3 fois. C'est là que la mémoire et l'oreille travaillent. \tabularnewline
7 & S'enregistrer et comparer & 1 fois. Écoutez votre version puis l'original. Notez \textbf{un} défaut à corriger demain. \tabularnewline
\end{tabularx}
\end{tableau}

\begin{astuce}[Variantes utiles]
\textbf{Vitesse réduite} (0,75×) les deux premières semaines, puis vitesse normale, puis
1,1× pour l'agilité. \textbf{Shadowing en marchant} : marcher en shadowant active le rythme
corporel et améliore nettement la prosodie. \textbf{Shadowing chuchoté} dans le bus : vous
travaillez sans gêner personne.
\end{astuce}

\begin{piege}[Ne shadowez pas votre propre accent]
Le but n'est pas de « dire le texte », mais de \textbf{copier le son}. Si vous lisez avec
votre accent français habituel, l'exercice ne sert presque à rien. Exagérez l'imitation,
comme un acteur : c'est la seule façon de casser vos automatismes.
\end{piege}

\section{Technique 2 --- La répétition espacée (Anki)}
\textbf{Ce que c'est :} un logiciel qui vous montre chaque carte au moment exact où vous
alliez l'oublier. C'est le multiplicateur de vocabulaire.

\subsection{Régler Anki en 5 minutes}
\begin{enumerate}
  \item Créez un paquet nommé \en{Anglais --- 12 mois}.
  \item Type de carte : \textbf{Phrase à trou} (\en{Cloze}) de préférence --- bien plus
  efficace que « mot isolé ».
  \item Cartes nouvelles par jour : \textbf{15} (au début), jusqu'à 25 à partir du mois 3.
  \item Intervalle maximal : 180 jours. Pas de limite de révisions par jour.
  \item Activez l'audio automatique si disponible (ou ajoutez le son vous-même depuis Forvo).
\end{enumerate}

\subsection{Comment fabriquer une bonne carte}
\begin{tableau}[Mauvaise carte \ra bonne carte]
\begin{tabularx}{\linewidth}{@{}p{7.6cm}X@{}}
\entete{\ko À éviter} & \entete{\ok À faire} \tabularnewline
\hline
Recto : \en{to achieve} \quad Verso : \en{réussir, accomplir} &
Recto : \en{She finally \{\{c1::achieved\}\} her goal.} \\
& Verso : \en{achieve} = atteindre, accomplir --- \pron{a-tchi-v} \\
& + la phrase entière lue à voix haute. \tabularnewline
\hline
Un mot seul, hors contexte & Toujours \textbf{une phrase complète}, personnelle si possible (votre vie, votre travail). \tabularnewline
\hline
Une carte = 5 sens possibles & \textbf{Une seule réponse attendue} par carte. \tabularnewline
\hline
Traduire FR \ra EN uniquement & Les deux sens : EN \ra FR (compréhension) et FR \ra EN (production). \tabularnewline
\end{tabularx}
\end{tableau}

\begin{memo}[Les 3 règles d'or d'Anki]
\textbf{1.} Faites vos révisions \textbf{tous les jours}, à la même heure (15 min suffisent).
\textbf{2.} Dites chaque carte \textbf{à voix haute}, en entier, avec l'intonation.
\textbf{3.} Ne créez que des cartes pour des mots que vous avez \textbf{réellement rencontrés}
(dans ce livre, un podcast, une série) --- jamais des listes apprises « au cas où ».
\end{memo}

\section{Technique 3 --- L'auto-conversation (\emph{self-talk})}
\textbf{Ce que c'est :} parler seul en anglais, à voix haute ou dans sa tête, en commentant
ce que vous faites. C'est gratuit, disponible 24 h/24, et cela entraîne exactement le muscle
qui vous manque : \textbf{produire sans interlocuteur pour vous aider}.

\begin{tableau}[Les 5 niveaux du self-talk --- progressez dans l'ordre]
\begin{tabularx}{\linewidth}{@{}p{1.4cm}p{3.6cm}X@{}}
\entete{Niveau} & \entete{Exercice} & \entete{Exemple} \tabularnewline
\hline
\textbf{1} & Nommer & Vous regardez autour de vous et nommez tout ce que vous voyez, à voix haute, 2 minutes. \en{a table, a phone, the window, my shoes\ldots} \tabularnewline
\textbf{2} & Commenter ses actes & \en{I'm opening the fridge. I'm looking for some water. It's cold.} (en cuisinant, en vous habillant, en marchant) \tabularnewline
\textbf{3} & Raconter sa journée & Chaque soir, 3 minutes : \en{This morning I woke up at six. Then I\ldots} \tabularnewline
\textbf{4} & Anticiper & Le matin : \en{Today I'm going to\ldots\ First I will\ldots\ Then\ldots} \tabularnewline
\textbf{5} & Argumenter / débatteur & Choisissez un sujet (\en{Is social media good for young people?}) et parlez 3 minutes \emph{pour}, puis 3 minutes \emph{contre}. \tabularnewline
\end{tabularx}
\end{tableau}

\begin{astuce}[Le « trou de mot » est votre meilleur professeur]
Quand un mot manque en self-talk, \textbf{ne cherchez pas tout de suite}. Contournez-le :
\en{``It's the thing you use to open a bottle\ldots''}. Puis notez le mot manquant sur votre
téléphone et créez une carte Anki le soir. Vous ne l'oublierez plus jamais, parce que votre
cerveau l'a \emph{cherché}.
\end{astuce}

\section{Technique 4 --- Le journal en anglais}
Dix minutes par jour, cinq à dix phrases. Aucun dictionnaire au premier jet --- écrivez avec
ce que vous avez, puis corrigez.

\begin{tableau}[Prompts de journal --- 30 sujets pour le premier mois]
\footnotesize
\begin{multicols}{2}
\begin{enumerate}
\item Ce que j'ai fait hier.
\item Ma famille.
\item Mon quartier / ma ville.
\item Mon travail ou mes études.
\item Ce que j'ai mangé aujourd'hui.
\item Mes projets pour cette année.
\item La personne que j'admire le plus.
\item Mon meilleur souvenir d'enfance.
\item Ce qui m'a rendu heureux cette semaine.
\item Ce qui m'a énervé cette semaine.
\item Ma journée idéale.
\item Le temps qu'il fait et ce que je fais avec.
\item Un film ou une série que j'ai vu.
\item Ce que je veux dans cinq ans.
\item Mon téléphone : à quoi il me sert.
\item Une habitude que je veux changer.
\item Un voyage que j'ai fait.
\item La description de ma chambre.
\item Ce que j'ai appris cette semaine en anglais.
\item Un problème que j'ai résolu.
\item Mon plat préféré, comment on le cuisine.
\item Une conversation que j'ai eue aujourd'hui.
\item Ce dont j'ai peur.
\item Mes trois objectifs du mois.
\item Une personne que je veux remercier.
\item Ce que je ferais avec un million.
\item Ma routine du matin, minute par minute.
\item Une chose nouvelle que j'ai essayée.
\item Ce qui est difficile en anglais pour moi.
\item Bilan de mon mois d'anglais.
\end{enumerate}
\end{multicols}
\end{tableau}

\section{Technique 5 --- La dictée (\en{dictation})}
Écrire ce que l'on entend : c'est le seul exercice qui prouve si vous entendez \emph{vraiment}.
Protocole (10 minutes) :
\begin{enumerate}
  \item Écoutez une phrase (5 à 10 mots) une fois, sans écrire.
  \item Écoutez-la à nouveau et écrivez ce que vous entendez, en mettant un tiret pour chaque
  trou : \en{I --- to the --- yesterday}.
  \item Réécoutez jusqu'à 5 fois en complétant.
  \item Comparez avec le texte exact. Comptez les trous.
  \item \textbf{Analysez chaque erreur} : était-ce un mot inconnu ? un mot connu mais « mangé »
  par la liaison (\en{going to} entendu comme \en{gonna}) ? un son non perçu (\en{th}, \en{h} final) ?
  \item Relisez la phrase à voix haute trois fois, en imitant l'audio.
\end{enumerate}
Cinq phrases par jour, cinq jours par semaine : au bout de trois mois, votre oreille n'a plus
rien à voir.

\section{Technique 6 --- La lecture à voix haute}
Prenez n'importe quel texte de ce livre (dialogue, exemple, paragraphe) et lisez-le à voix
haute, lentement, en articulant, 5 minutes par jour. Bénéfices : votre bouche s'habitue aux
enchaînements anglais ; vous mémorisez des structures par la voie motrice ; vous repérez
immédiatement les mots que vous « savez » mais ne savez pas dire.

\begin{astuce}[La règle des trois lectures]
\textbf{1\textsuperscript{re} lecture :} lente, mot à mot, en corrigeant la prononciation.
\textbf{2\textsuperscript{e} lecture :} normale, en respectant les pauses et l'accent tonique.
\textbf{3\textsuperscript{e} lecture :} rapide et expressive, comme un acteur --- en levant
les yeux du texte à la fin de chaque phrase.
\end{astuce}

\section{Technique 7 --- L'immersion fabriquée}
Vous ne vivez pas à Londres ? Fabriquez Londres chez vous. Voici la liste à mettre en place
\textbf{dans la première semaine}, une fois pour toutes :

\begin{tableau}[Votre environnement d'immersion --- à installer une fois]
\begin{tabularx}{\linewidth}{@{}p{4.4cm}X@{}}
\entete{Élément} & \entete{Action concrète} \tabularnewline
\hline
\textbf{Téléphone et ordinateur} & Langue du système en anglais. Vous apprendrez 200 mots utiles en trois jours sans effort. \tabularnewline
\textbf{Réseaux sociaux} & Abonnez-vous à 10 comptes anglophones dans des domaines que vous aimez déjà (football, cuisine, religion, technologie, humour, actualité). L'algorithme vous servira de l'anglais toute la journée. \tabularnewline
\textbf{Musique} & Playlist anglaise pour la cuisine, le transport, le sport. Une chanson travaillée par semaine (paroles + sens + chant). \tabularnewline
\textbf{Séries et films} & Sous-titres \textbf{anglais} (jamais français) à partir du mois 3. Avant : sous-titres français, mais 10 minutes par épisode en sous-titres anglais. \tabularnewline
\textbf{Radio / podcast} & Un flux anglais en fond sonore pendant les tâches ménagères, même sans comprendre. Le cerveau s'habitue au rythme. \tabularnewline
\textbf{YouTube} & Historique de recommandations à « nettoyer » : regardez 20 vidéos en anglais, l'algorithme basculera. \tabularnewline
\textbf{Étiquettes} & Collez des post-it en anglais sur 20 objets de la maison (\en{fridge, mirror, cupboard, kettle, shelf\ldots}). \tabularnewline
\textbf{Pensée} & Deux « alarmes » par jour (11 h et 17 h) : au bip, pensez en anglais pendant 2 minutes, quoi qu'il arrive. \tabularnewline
\textbf{Personnes} & Un partenaire de langue identifié (application, groupe, club, collègue, camarade de classe) --- rendez-vous hebdomadaire fixé dans l'agenda. \tabularnewline
\end{tabularx}
\end{tableau}

\section{Technique 8 --- L'enregistrement de soi}
Une fois par semaine, enregistrez 3 minutes de parole libre sur un sujet donné. Écoutez
le lendemain --- pas immédiatement (vous seriez trop sévère) --- avec une grille de trois
questions :
\begin{enumerate}
  \item \textbf{Fluidité :} combien d'hésitations de plus de 3 secondes ? combien de retours
  au français ?
  \item \textbf{Prononciation :} quels sons ont dérapé (\en{th}, \en{h}, voyelles, accent tonique) ?
  \item \textbf{Grammaire :} quelles erreurs reviennent ? (Notez-les : trois maximum, pas plus.
  Travaillez-les la semaine suivante.)
\end{enumerate}

\begin{memo}[Le cycle hebdomadaire complet]
\textbf{Lundi :} nouvelle unité. \textbf{Mardi :} révision + shadowing.
\textbf{Mercredi :} conversation avec partenaire. \textbf{Jeudi :} lecture + grammaire.
\textbf{Vendredi :} dictée + chanson. \textbf{Samedi :} projet (série, présentation).
\textbf{Dimanche :} enregistrement de 3 min + test rapide + bilan écrit.
\end{memo}


% =====================================================================
\chapter{Outils, ressources et partenaires}

Tout ce dont vous avez besoin est \textbf{gratuit} ou presque. Voici l'équipement minimal,
puis les ressources classées par niveau. (Les liens détaillés sont en annexe E.)

\section{L'équipement minimal (à réunir en une soirée)}
\begin{enumerate}
  \item \textbf{Un téléphone} avec un enregistreur vocal et des écouteurs. C'est l'outil
  central : il remplace un laboratoire de langue.
  \item \textbf{Anki} (ordinateur ou téléphone) --- répétition espacée. Gratuit sur Android
  (\en{AnkiDroid}) et ordinateur ; payant sur iPhone (\en{AnkiMobile}), où \en{AnkiWeb}
  gratuit suffit.
  \item \textbf{Un dictionnaire anglais-anglais} : Cambridge Dictionary ou Oxford Learner's
  Dictionaries. Passer au dictionnaire unilingue dès le mois 3 est un accélérateur majeur.
  \item \textbf{Un dictionnaire bilingue de secours} : WordReference (avec ses forums, très
  utiles pour les nuances) ou Linguee/DeepL (pour voir des traductions en contexte réel).
  \item \textbf{Un correcteur} : LanguageTool ou le correcteur de Google Docs, pour vos
  journaux écrits. Attention : lisez la correction, ne la copiez pas bêtement.
  \item \textbf{Ce livre}, imprimé ou sur tablette, plus un cahier de 100 pages.
\end{enumerate}

\section{Ressources d'écoute classées par niveau}
\begin{tableau}[Où écouter, selon votre mois]
\begin{tabularx}{\linewidth}{@{}p{2.0cm}p{4.6cm}X@{}}
\entete{Niveau} & \entete{Type de ressource} & \entete{Exemples et usage} \tabularnewline
\hline
\textbf{A1} (mois 1-2) & Podcasts pour apprenants, débit très lent &
Chaînes et podcasts de type « English for beginners », \en{BBC Learning English}
(rubriques débutants), \en{VOA Learning English} (débit lent, vocabulaire contrôlé).
15 min/jour, vitesse 0,75×. \tabularnewline
\textbf{A2} (mois 3-5) & Histoires courtes, dialogues lents, chansons claires &
Podcasts d'histoires pour apprenants (\en{storytelling} lent), \en{6 Minute English} (BBC),
dessins animés pour enfants (vocabulaire simple, articulation exagérée --- ne les méprisez
pas : c'est un excellent matériau). \tabularnewline
\textbf{B1} (mois 6-8) & Podcasts « naturels mais articulés », TED-Ed, séries &
Conférences TED courtes sous-titrées, documentaires, sitcoms (\en{Friends},
\en{Modern Family}), chaînes YouTube de vulgarisation. \tabularnewline
\textbf{B2} (mois 9-12) & Contenu natif : actualités, interviews, humour &
Journaux télévisés (BBC, CNN, Al Jazeera English), podcasts natifs sur vos centres d'intérêt,
débats, stand-up, films sans sous-titres. \tabularnewline
\textbf{Accents} (mois 11) & Variétés d'anglais &
Cherchez volontairement un accent britannique, un américain, un nigérian, un ghanéen, un
kényan, un sud-africain, un indien. Votre oreille doit être mondiale, pas seulement
hollywoodienne. \tabularnewline
\end{tabularx}
\end{tableau}

\section{Trouver des gens à qui parler (le point le plus important du chapitre)}
\begin{tableau}[Huit façons de parler à un humain, de la plus facile à la plus efficace]
\begin{tabularx}{\linewidth}{@{}p{4.6cm}X@{}}
\entete{Solution} & \entete{Comment faire} \tabularnewline
\hline
\textbf{Applications d'échange linguistique} &
HelloTalk, Tandem : vous aidez quelqu'un en français/wolof, il vous aide en anglais.
Messages vocaux d'abord, appels ensuite. 20 min/jour possible. \tabularnewline
\textbf{Serviteurs Discord / groupes WhatsApp} &
De nombreux serveurs d'apprentissage de l'anglais ont des salons vocaux ouverts 24 h/24.
Cherchez « English learning voice chat ». Gratuit, immédiat, anonyme --- donc sans honte. \tabularnewline
\textbf{Clubs de conversation locaux} &
Instituts culturels, centres américains/britanniques, bibliothèques, universités, cafés
langues. Une séance par semaine, souvent gratuite. \tabularnewline
\textbf{Communautés religieuses} &
Dans beaucoup de villes d'Afrique de l'Ouest, des églises, des groupes d'étude ou des
associations tiennent des cultes ou séances en anglais. Renseignez-vous : c'est une source
régulière et bienveillante de pratique. \tabularnewline
\textbf{Un camarade de même niveau} &
La solution la plus sous-estimée : trouvez une personne de votre entourage qui veut aussi
apprendre, et imposez-vous \textbf{une heure par semaine en anglais uniquement}, même
imparfait. Vous progresserez plus vite qu'avec un professeur que vous voyez une fois par mois. \tabularnewline
\textbf{Professeur particulier en ligne} &
Plateformes de tutorat (payantes) : 30 min, une à deux fois par semaine. Utile surtout pour
corriger la prononciation et débloquer l'oral à partir du mois 4. \tabularnewline
\textbf{Parler à un assistant vocal} &
L'assistant de votre téléphone, réglé en anglais, comprend et répond. Excellent pour
s'entraîner sans jugement : \en{``What's the weather today?'', ``Set a timer for ten
minutes'', ``Tell me a joke.''} \tabularnewline
\textbf{Parler seul (self-talk)} &
Toujours disponible, jamais en panne, jamais moqueur. C'est 80\,\% du volume de parole
de la plupart des autodidactes qui réussissent. \tabularnewline
\end{tabularx}
\end{tableau}

\begin{piege}[La honte est votre seul véritable ennemi]
Les francophones ne manquent ni d'intelligence ni de ressources : ils manquent
d'\textbf{audace}. Vous ferez des fautes. Vous aurez l'accent. Vous chercherez vos mots.
Les anglophones s'en moquent --- ils sont même admiratifs, parce qu'eux ne parlent qu'une
langue. Répétez-vous cette phrase chaque fois que vous hésitez à ouvrir la bouche :
\en{``Done is better than perfect.''} \textit{(Fait vaut mieux que parfait.)}
\end{piege}

\section{Trois routines à copier}

\subsection{Routine « trajet » (20 min, deux fois par jour)}
\begin{enumerate}
  \item 5 min : révisions Anki (téléphone en main).
  \item 10 min : shadowing d'un épisode de podcast déjà écouté.
  \item 5 min : self-talk --- raconter à voix basse ce que vous allez faire aujourd'hui.
\end{enumerate}

\subsection{Routine « soir calme » (45 min)}
\begin{enumerate}
  \item 15 min : unité de grammaire + exercices.
  \item 15 min : un épisode de série, 3 passes (voir chapitre 16).
  \item 10 min : journal écrit.
  \item 5 min : relecture à voix haute de 10 phrases de la journée.
\end{enumerate}

\subsection{Routine « dimanche bilan » (60 min)}
\begin{enumerate}
  \item 10 min : test rapide de la semaine (questions du livre).
  \item 15 min : enregistrement de 3 min + écoute + 3 défauts notés.
  \item 15 min : nettoyage des cartes Anki (supprimer les cartes inutiles, corriger les
  mauvaises).
  \item 10 min : choix du matériel de la semaine prochaine.
  \item 10 min : conversation avec un partenaire (même 10 minutes comptent).
\end{enumerate}

\begin{mission}[Semaine 1 --- installer la machine]
\begin{enumerate}
  \item Jour 1 : téléphone en anglais + Anki installé + 20 post-it collés.
  \item Jour 2 : premier shadowing (15 min) sur un audio de 1 minute.
  \item Jour 3 : premier journal (5 phrases) + première dictée (3 phrases).
  \item Jour 4 : auto-conversation niveau 1 et 2 (5 min).
  \item Jour 5 : trouver et contacter un partenaire de langue (un seul message suffit).
  \item Jour 6 : enregistrement de 2 minutes sur « pourquoi j'apprends l'anglais ».
  \item Jour 7 : bilan + cases cochées + récompense (vous en aurez besoin toute l'année :
  prévoyez une petite récompense chaque dimanche).
\end{enumerate}
\end{mission}
